% ============================================================
% CHAPTER 5 --- PREDICATE LOGIC  (CN Ch.5, §5.1-5.15)
% ============================================================
\newpage
\section{Predicate Logic}

\subsection{Relations, predicates, and truth-valued functions}
\label{sec:5-1}
% CN 5.1 / Rosen 1.4

\begin{definition}[Predicate]
A \emph{predicate} is a relation (\S\ref{sec:2-13}, \S\ref{sec:2-17}), used with intent to do logic with it. A predicate with $k$ arguments is \emph{$k$-ary}.
\end{definition}

\begin{definition}[Predicate as a truth-valued function]
A $k$-ary relation $R \subseteq A_1\times\cdots\times A_k$ can equivalently be viewed as a function
\[
R : A_1\times\cdots\times A_k \to \{T,F\}, \qquad R(a_1,\dots,a_k) = T \iff (a_1,\dots,a_k)\in R.
\]
Rosen also calls this a \emph{propositional function}.
\end{definition}

\begin{note}[Substitution yields a proposition]
Giving specific values to every argument of a predicate produces a proposition (definite truth value).
\end{note}

\begin{note}[Equality is always available]
$=$ is a predicate assumed available for \emph{every} domain -- without it, no domain of objects could even be meaningfully discussed.
\end{note}

\begin{definition}[Property]
A unary predicate (\emph{property}) $P(x)$ on domain $D$ corresponds to the set $\{x\in D : P(x)\}$.
\end{definition}

\subsection{Variables and constants}
\label{sec:5-2}
% CN 5.2 / Rosen 1.4

\begin{definition}[Variable]
A name for an object, without specifying it exactly, used to reason about an unknown or arbitrary member of its \emph{domain}.
\end{definition}

\begin{note}[Domain]
    $Domain$ also call $domain of discourse$ or $universe of discourse$.
\end{note}

\begin{definition}[Constant]
An object belonging to the domain of some variable -- a specific, named member of that domain.
\end{definition}

\begin{note}[Every variable has a domain]
E.g.\ ``let $x\in\Z$'' declares $x$ with domain $\Z$; a Boolean variable always has domain $\{T,F\}$.
\end{note}

\subsection{Predicates and variables}
\label{sec:5-3}
% CN 5.3 / Rosen 1.4

\begin{definition}[Free variable]
A variable in a predicate expression not (yet) assigned a value.
\end{definition}

\begin{note}[Truth value requires all arguments fixed]
An expression with $\ge1$ free variable is not a proposition (no truth value). Only when \emph{every} free variable is assigned a value from its domain does the expression become a proposition.
\end{note}

\begin{note}[Domain matching is mandatory]
Any value or variable placed into an argument of a predicate must belong to that argument's declared domain -- mismatched domains make the expression ill-formed, not just false.
\end{note}

\begin{note}[Prefix vs.\ infix notation]
Newly defined predicates typically use prefix notation, $R(a,b)$; predicates for ordering/equality/equivalence conventionally use infix notation, $a<b$, $a=b$ (\S\ref{sec:2-5}).
\end{note}

\subsection{Arguments of predicates}
\label{sec:5-4}
% CN 5.4

\begin{definition}[Term, recursive]
A \emph{term} of domain $D$ is one of:
\begin{itemize}
    \item a constant from $D$;
    \item a variable with domain $D$;
    \item a function with codomain $D$, each of whose arguments is given a term of matching domain.
\end{itemize}
\end{definition}

\begin{note}[Recursion bottoms out]
Terms are built from constants/variables/functions to arbitrary finite depth (composition, nested function calls), but never infinitely -- every term has a finite construction.
\end{note}

\begin{example}
Let $a,b$ be variables with domain $\R$. The expression
\[
a^2+b^2
\]
is a term of domain $\R$.

To see this, first consider $a^2$. The squaring function
\[
\operatorname{square}:\R\to\R
\]
has codomain $\R$, and its argument $a$ is a term of domain $\R$. Therefore,
\[
a^2=\operatorname{square}(a)
\]
is a term of domain $\R$.

Similarly, since $b$ is a term of domain $\R$,
\[
b^2=\operatorname{square}(b)
\]
is also a term of domain $\R$.

Finally, consider the real addition function
\[
+:\R\times\R\to\R.
\]
Its two arguments must each be terms of domain $\R$. Since both $a^2$ and $b^2$ are terms of domain $\R$, they can be given to the two arguments of the addition function. Hence,
\[
a^2+b^2
\]
is a term of domain $\R$.
\end{example}

\begin{example}
Let $\mathbb P$ be the domain of people, and suppose that
\[
\operatorname{FatherOf}:\mathbb P\to\mathbb P
\]
and
\[
\operatorname{MotherOf}:\mathbb P\to\mathbb P.
\]
If $Y$ is a variable with domain $\mathbb P$, then
\[
\operatorname{FatherOf}(Y)
\]
is a term of domain $\mathbb P$. Consequently,
\[
\operatorname{MotherOf}(\operatorname{FatherOf}(Y))
\]
is also a term of domain $\mathbb P$.
\end{example}

\subsection{Building logical expressions with predicates}
\label{sec:5-5}
% CN 5.5 / Rosen 1.4

\begin{definition}[Predicate logic expression, recursive]
\begin{itemize}
    \item $T, F$ are predicate logic expressions.
    \item A predicate applied to terms of matching domains is a predicate logic expression.
    \item If $E,F$ are predicate logic expressions, so are $(E),\ \neg E,\ E\wedge F,\ E\vee F,\ E\implies F,\ E\iff F$.
\end{itemize}
\end{definition}

\begin{note}[Free variables propagate]
A predicate logic expression's free variables are the union of the free variables of its terms/sub-expressions -- combining expressions never removes a free variable (only quantification, \S\ref{sec:5-6}-5.9, does that).
\end{note}

\begin{example}
$(1<i)\wedge(i<n)$ has free variables $i,n$. $\mathrm{Parent}(X,Y)\iff\big(Y=\mathrm{MotherOf}(X)\big)\vee\big(Y=\mathrm{FatherOf}(X)\big)$ has free variables $X,Y$.
\end{example}

\subsection{Existential quantifier}
\label{sec:5-6}
% CN 5.6 / Rosen 1.4

\begin{definition}[Existential quantifier]
For predicate $P(x)$ on domain $D$:
\[
\exists x\, P(x)
\]
is true iff $P(a)$ is true for \emph{at least one} $a\in D$.
\end{definition}

\begin{note}[Punctuation]
$\exists x : P(x)$, $\exists x.\, P(x)$, $\exists x\, P(x)$ are all used; the colon/period is read ``such that''.
\end{note}

\begin{definition}[Bound variable]
Once quantified, $x$ is \emph{bound}: no longer available to receive a specific value. A fully-quantified expression has no free variables and is a genuine proposition, with a definite truth value.
\end{definition}

\begin{note}[Domain matters]
$\exists W\in\N:\ W<0$ is false; $\exists W\in\Z:\ W<0$ is true -- same predicate, different domain, different truth value.
\end{note}

\begin{note}[Analogy with disjunction -- and where it breaks]
$\exists x \in D,\ P(x)$ behaves like the disjunction $\bigvee_{a\in D} P(a)$ -- true iff at least one disjunct is true. But this is only an analogy: disjunctions must be \emph{finite}, while $D$ may be infinite, so existential quantification genuinely extends what propositional logic alone can express.
\end{note}

\begin{note}[Vacuous quantification]
If $x$ does not appear in $P$, $\exists x\, P$ is equivalent to $P$ itself (the quantifier is inert): e.g.\ $\exists w:\ z<0$ is just $z<0$.
\end{note}

\begin{note}
Quantifiers apply only to variables, never to constants: $\exists 5$ or $\exists\text{Annie}$ are ill-formed.
\end{note}

\subsection{Restricting existentially quantified variables}
\label{sec:5-7}
% CN 5.7 / Rosen 1.4 (restricted quantifiers)

\begin{note}[Restricting $\exists$ uses $\wedge$]
To assert ``some $X$ satisfying $R(X)$ also satisfies $P(X)$'':
\[
\exists X:\ R(X) \wedge P(X).
\]
\end{note}

\subsection{Universal quantifier}
\label{sec:5-8}
% CN 5.8 / Rosen 1.4

\begin{definition}[Universal quantifier]
For predicate $P(x)$ on domain $D$:
\[
\forall x\, P(x)
\]
is true iff $P(a)$ is true for \emph{every} $a\in D$.
\end{definition}

\begin{note}
As with $\exists$: the quantified variable becomes bound; an irrelevant quantifier (variable not appearing in the body) can be dropped.
\end{note}

\bigskip
\begin{note}[Vacuous truth generalizes]
$\forall x\in\emptyset,\ P(x)$ is vacuously true for \emph{any} $P$, exactly generalizing $\emptyset\subseteq A$ (\S\ref{sec:1-6} Theorem 1) to arbitrary predicates, not just set membership.
\end{note}

\subsection{Restricting universally quantified variables}
\label{sec:5-9}
% CN 5.9 / Rosen 1.4

\begin{note}[Restricting $\forall$ uses $\implies$]
To assert ``every $X$ satisfying $R(X)$ also satisfies $P(X)$'':
\[
\forall X:\ R(X) \implies P(X).
\]
\end{note}

\bigskip
\begin{note}[Preview: De Morgan's laws for quantifiers]
\[
\neg\forall x\, P(x) \ \equiv\ \exists x\,\neg P(x), \qquad \neg\exists x\, P(x) \ \equiv\ \forall x\,\neg P(x).
\]
\end{note}
\begin{note}
Consequence: negating $\forall X: R(X)\implies P(X)$ gives $\exists X: R(X)\wedge\neg P(X)$ -- exactly matching the $\exists/\wedge$ pairing above, since $\neg(R\implies P) \equiv R\wedge\neg P$ (\S\ref{sec:4-15}).
\end{note}

\subsection{Multiple quantifiers}
\label{sec:5-10}
% CN 5.10 / Rosen 1.4

\begin{note}[Same-type quantifiers commute]
\[
\exists X\,\exists Y\, P = \exists Y\,\exists X\, P = \exists(X,Y)\, P, \qquad
\forall X\,\forall Y\, P = \forall Y\,\forall X\, P = \forall(X,Y)\, P.
\]
\end{note}

\begin{note}[Mixed quantifiers do \emph{not} commute]
\[
\exists X\,\forall Y\, P, \quad \forall X\,\exists Y\, P, \quad \exists Y\,\forall X\, P, \quad \forall Y\,\exists X\, P
\]
are, in general, four genuinely different statements.
\end{note}

\bigskip
\begin{example}[Classic order-matters example]
Let $D(x,y)$ mean ``$x$ has dated $y$'' on the domain of all people.
\[
\exists x\,\forall y\, D(x,y): \text{someone has dated everyone}.
\]
\[
\forall y\,\exists x\, D(x,y): \text{everyone has been dated by someone.}
\]
The first implies the second (that one exceptional dater witnesses every $y$), but not conversely -- the second only needs, for each $y$, \emph{some} $x$, possibly different for each $y$.
\end{example}

\begin{proposition}[$\exists\forall \implies \forall\exists$, general fact]
\[
\exists x\,\forall y\, P(x,y) \ \implies\ \forall y\,\exists x\, P(x,y).
\]
The converse implication does \emph{not} hold in general.
\end{proposition}
\begin{proof}
Suppose $\exists x\,\forall y\,P(x,y)$: fix such an $x_0$, so $P(x_0,y)$ holds for every $y$. Then for any $y$, taking $x=x_0$ witnesses $\exists x\,P(x,y)$ -- so $\forall y\,\exists x\,P(x,y)$ holds. The dating example above (with $D$ not symmetric enough to reverse) shows the converse can fail.
\end{proof}

\paragraph{Quantifier--connective pitfall.***} $\forall$ pairs with
$\Rightarrow$; $\exists$ pairs with $\wedge$ --- the other way round
lets a trivial witness make the statement true for the wrong reason.

\subsection{Predicate logic expressions}
\label{sec:5-11}
% CN 5.11 / Rosen 1.4

\begin{definition}[Predicate logic expression, complete recursive grammar]
\begin{itemize}
    \item $T,F$ are predicate logic expressions.
    \item Terms of matching domains plugged into a predicate give a predicate logic expression.
    \item If $E,F$ are predicate logic expressions: $(E),\ \neg E,\ E\wedge F,\ E\vee F,\ E\implies F,\ E\iff F$.
    \item If $E$ is a predicate logic expression and $X$ is a free variable of $E$: $\exists X: E$ and $\forall X: E$ are predicate logic expressions, and $X$ is \emph{no longer free} in either.
\end{itemize}
\end{definition}

\begin{note}[Free variables shrink under quantification]
If $V$ is the free-variable set of $E$, then $\exists X: E$ and $\forall X: E$ each have free-variable set $V\setminus\{X\}$.
\end{note}

\subsection{Doing logic with quantifiers}
\label{sec:5-12}
% CN 5.12 / Rosen 1.4

\begin{note}[Instantiation and generalization]
\[
\big(\forall X\, P(X)\big) \implies P(\mathrm{obj}), \qquad P(\mathrm{obj}) \implies \big(\exists X\, P(X)\big)
\]
for any specific $\mathrm{obj}$ in $X$'s domain.
\end{note}

\begin{proposition}[$\forall$ distributes over $\wedge$; $\exists$ distributes over $\vee$]
\[
\forall X\,(P(X)\wedge Q(X)) \ \equiv\ (\forall X\, P(X)) \wedge (\forall X\, Q(X)),
\]
\[
\exists X\,(P(X)\vee Q(X)) \ \equiv\ (\exists X\, P(X)) \vee (\exists X\, Q(X)).
\]
\end{proposition}

\begin{note}[$\forall$ over $\vee$, and $\exists$ over $\wedge$, only give one-way implications]
\[
(\forall X\, P(X)) \vee (\forall X\, Q(X)) \implies \forall X\,(P(X)\vee Q(X)),
\]
\[
\exists X\,(P(X)\wedge Q(X)) \implies (\exists X\, P(X)) \wedge (\exists X\, Q(X)),
\]
but neither converse holds in general -- e.g.\ every integer is even or odd ($\forall X$, $\mathrm{even}(X)\vee\mathrm{odd}(X)$), yet it is false that every integer is even, or every integer is odd.
\end{note}

\begin{definition}[Scope of a quantified variable]
The sub-expression from the quantifier to the end of its innermost enclosing parentheses (or the end of the whole expression, if unparenthesized).
\end{definition}

\begin{example}[Two disjoint scopes, same variable name]
In $(\forall X\, P(X)) \wedge (\forall X\, Q(X))$, the two $X$'s are entirely independent -- like distinct local variables in a program that happen to share a name.
\end{example}

\paragraph{At least / at most / exactly $n$ (predicate logic).}
At least $n$:
\[
\exists x_1\cdots x_n\Big(\bigwedge_{i<j}x_i\neq x_j \wedge \bigwedge_i P(x_i)\Big)
\]
At most $n$:
\[
\forall x_1\cdots x_{n+1}\Big(\bigwedge_i P(x_i)\Rightarrow \bigvee_{i<j}x_i=x_j\Big)
\]
Exactly $n$: conjunction of both.

\paragraph{The uniqueness quantifier $\exists!$.} $\exists!\,x\,P(x)$
means ``there is a unique $x$ such that $P(x)$.'' Without ``$!$'':
\[
\exists!\,x\,P(x) \;\equiv\; \exists x\,\Big(P(x)\wedge\forall z\big(P(z)\Rightarrow z=x\big)\Big)
\]
General template for ``exactly one thing satisfies $P$'':
existence ($P(x)$) $\wedge$ uniqueness (anything satisfying $P$
equals $x$). \emph{(CN Ch.5 Problem 4.)}

\subsection{Duality between quantifiers}
\label{sec:5-13}
% CN 5.13 / Rosen 1.4 (De Morgan's laws for quantifiers)

\begin{theorem}[De Morgan's laws for quantifiers]
\[
\neg\forall X\, P(X) \ \equiv\ \exists X\,\neg P(X), \qquad \neg\exists X\, P(X) \ \equiv\ \forall X\,\neg P(X).
\]
\end{theorem}

\begin{example}[``Not all stars are visible'' $=$ ``some star is invisible'']
\begin{align*}
\neg\forall X\,(\mathrm{star}(X)\implies\mathrm{visible}(X))
&= \exists X\,\neg(\mathrm{star}(X)\implies\mathrm{visible}(X)) \\
&= \exists X\,\neg(\neg\mathrm{star}(X)\vee\mathrm{visible}(X)) \\
&= \exists X\,(\mathrm{star}(X)\wedge\neg\mathrm{visible}(X)).
\end{align*}
\end{example}

\begin{note}
$\neg\forall Y\,\neg = \exists Y$, and $\neg\exists Y\,\neg = \forall Y$ (apply the laws twice, using double negation).
\end{note}

\subsection{Some examples}
\label{sec:5-15}
% CN 5.15 / Rosen 1.4

\begin{example}[``Infinitely many primes'', in predicate logic]
With $\mathrm{Prime}(n)$ on domain $\N$ and predicate $\le$:
\[
\forall b\, \exists n\ \mathrm{Prime}(n) \wedge \neg(n \le b).
\]
Read: for every bound $b$, some prime exceeds it -- correctly capturing ``infinitely many primes'' without needing an infinite conjunction.
\end{example}

\begin{example}[Swapping the quantifiers gives a different, false, statement]
\[
\exists n\, \forall b\ \mathrm{Prime}(n) \wedge \neg(n \le b).
\]
\end{example}
\begin{proof}[Why this is false -- step by step]
Distribute $\forall b$ over the conjunction (\S\ref{sec:5-12}):
\[
\forall b\ \mathrm{Prime}(n) \wedge \neg(n\le b) \ \equiv\ (\forall b\, \mathrm{Prime}(n)) \wedge (\forall b\, \neg(n\le b)).
\]

Since $\mathrm{Prime}(n)$ does not involve $b$, the quantifier is superfluous (\S\ref{sec:5-9}):
\[
\forall b\, \mathrm{Prime}(n) \ \equiv\ \mathrm{Prime}(n).
\]

Substituting back:
\[
\exists n\,\big(\mathrm{Prime}(n) \wedge \forall b\,\neg(n\le b)\big) \ \equiv\ \exists n\in\{\text{primes}\}\ \forall b\ n>b.
\]

This asserts a single prime exceeding \emph{every} number -- false. So the two quantifier orderings genuinely differ in meaning, one true and one false.
\end{proof}

\bigskip
\begin{example}[Translating a nested-quantifier English sentence]
``Every real number except $0$ has a multiplicative inverse.'' With $\mathrm{domain}=\R$:
\[
\forall x\,\big(x\neq 0 \implies \exists y\, (xy=1)\big).
\]
Note the restricted-$\forall$ pattern from \S\ref{sec:5-9} ($x\neq0 \implies \cdots$) nested with a plain $\exists y$ inside -- $y$'s existence is allowed to depend on $x$, which is essential: no single $y$ works for every nonzero $x$.
\end{example}

\subsection{Summary of rules for logic with quantifiers}
\label{sec:5-14}
% CN 5.14

\begin{note}[Full rule list]
\begin{align*}
\exists x\in A\, P(x) &\equiv \exists x\,(x\in A \wedge P(x)) & \text{(\S\ref{sec:5-7})}\\
\forall x\in A\, P(x) &\equiv \forall x\,(x\in A \implies P(x)) & \text{(\S\ref{sec:5-9})}\\
\exists x\, R &\equiv R, \quad \forall x\, R \equiv R & (x \text{ not free in } R;\ \S\ref{sec:5-7},5.9)\\
\exists x\exists y\, P(x,y) &\equiv \exists y\exists x\, P(x,y) & \text{(\S\ref{sec:5-10})}\\
\forall x\forall y\, P(x,y) &\equiv \forall y\forall x\, P(x,y) & \text{(\S\ref{sec:5-10})}\\
\forall x\, P(x) &\implies P(\mathrm{obj}), \quad P(\mathrm{obj}) \implies \exists x\, P(x) & \text{(\S\ref{sec:5-12})}\\
\exists x\,(P(x)\vee Q(x)) &\equiv (\exists x\,P(x))\vee(\exists x\,Q(x)) & \text{(\S\ref{sec:5-12})}\\
\forall x\,(P(x)\wedge Q(x)) &\equiv (\forall x\,P(x))\wedge(\forall x\,Q(x)) & \text{(\S\ref{sec:5-12})}\\
\neg\exists x\, P(x) &\equiv \forall x\,\neg P(x), \quad \neg\forall x\, P(x) \equiv \exists x\,\neg P(x) & \text{(\S\ref{sec:5-13})}
\end{align*}
\end{note}

